\chapter{ДИЗАЙН И ПРОВЕДЕНИЕ ЭКСПЕРИМЕНТОВ}

% -------------------------------------------------------
\section{Обзор экспериментального фреймворка}
% -------------------------------------------------------

\subsection{Три парадигмы как единый спектр}

Настоящее исследование объединяет три классические психологические парадигмы в едином экспериментальном фреймворке, обеспечивающем сопоставимость результатов на уровне моделей, метрик и условий доступа. Выбор парадигм обусловлен концептуальной логикой: они образуют прогрессирующий спектр сложности рассуждения и типа когнитивного предубеждения.

\textit{Эксперимент 1 (ошибка конъюнкции)} тестирует вероятностное суждение в одношаговом формате: модели предъявляется описание персонажа и два варианта ответа, один из которых нормативно правильный. Ошибка возникает в момент выбора и не требует многошагового вывода.

\textit{Эксперимент 2 (задача выбора карточек Уэйсона)} тестирует дедуктивное рассуждение: модель должна применить правило modus tollens к условному утверждению и выбрать карточки, способные его фальсифицировать. Задача является одношаговой, однако требует явного логического анализа структуры правила.

\textit{Эксперимент 3 (задача открытия правила 2-4-6)} тестирует индуктивное рассуждение в многоходовом интерактивном формате: модель должна итеративно формировать и пересматривать гипотезы на основе бинарной обратной связи. Предубеждение проявляется не в единственном выборе, а в стратегии тестирования на протяжении всего диалога.

Все три эксперимента проводились на одном и том же наборе из девяти моделей, что позволяет сравнивать профили предубеждений между парадигмами на уровне отдельных систем.

\subsection{Общие принципы дизайна}

Три принципа обеспечивают методологическую связность фреймворка.

\textit{Контроль контаминации.} В каждом эксперименте используются оригинальные датасеты, не воспроизводящие текст классических задач или их известных адаптаций. Детали контроля контаминации описаны в соответствующих подразделах.

\textit{Воспроизводимость.} Все модели запрашивались через единый API-интерфейс (OpenRouter). Сырые ответы сохранялись отдельно от производных метрик. Случаи нарушения формата фиксировались и обрабатывались через retry-логику.

\textit{Сопоставимость метрик.} Наряду со специфическими для каждого эксперимента показателями вводится единая метрика Severity of Bias, обеспечивающая сквозное сравнение интенсивности предубеждений между парадигмами.

